\section{Introdução}

Denomina-se problema de valor de contorno, ou PVC, o conjunto definido por uma equação diferencial parcial, ou simplesmente EDP, e suas condições de contorno, i.e. condições que delimitam o comportamento da solução na fronteira do domínio sobre o qual a equação está definida. Notoriamente, para EDPs nas quais há dependência temporal, faz-se necessário especificar, adicionalmente, a condição inicial; a este último caso, atribui-se de forma mais frequente o título de problema de valor inicial e de contorno \cite{Oliveira2010}. 

Esta classe de problemas, por assim dizer, permeia as mais diversas áreas da engenharia e das ciências exatas. Neste cenário, o problema é, em geral, definido não mais por uma única equação governante, mas sim por uma combinação de equações. Assim, no que tange ao estudo da dinâmica dos fluidos em meios porosos, por exemplo, o problema é governado por uma lei de conservação, a qual é derivada de princípios universais, tais como a conservação da massa, da energia e quantidade de movimento e, não obstante por uma lei constitutiva, ou seja, formulada através da modelagem física e baseada em observações experimentais \cite{Kirchdoerfer2016}.  

\textcolor{red}{Antes do início da bibliografia, deve-se contextualizar a temática denominada de \textit{data-drive computing}.}